\CatchFileBetweenTags{\AlphaInvVal}{constants.tex}{AlphaInvVal}

\CatchFileBetweenTags{\WeakAngleStepOneVal}{constants.tex}{WeakAngleStepOneVal}
\CatchFileBetweenTags{\WeakAngleVal}{constants.tex}{WeakAngleVal}
\CatchFileBetweenTags{\WBosonMassVal}{constants.tex}{WBosonMassVal}
\CatchFileBetweenTags{\AlphaInvTOL}{constants.tex}{AlphaInvTOL}
\CatchFileBetweenTags{\WeakAngleTCheckVal}{constants.tex}{WeakAngleTCheckVal}
\CatchFileBetweenTags{\WeakAngleExperimentalValue}{constants.tex}{WeakAngleExperimentalValue}
\CatchFileBetweenTags{\WeakAngleGlobalExperimentalValue}{constants.tex}{WeakAngleGlobalExperimentalValue}
\CatchFileBetweenTags{\WeakAngleGlobalSigma}{constants.tex}{WeakAngleGlobalSigma}

\CatchFileBetweenTags{\NcVal}{constants.tex}{NcVal}
\CatchFileBetweenTags{\NcEq}{constants.tex}{NcEq}
\CatchFileBetweenTags{\WeinbergGUTVal}{constants.tex}{WeinbergGUTVal}
\CatchFileBetweenTags{\WeinbergGUTEq}{constants.tex}{WeinbergGUTEq}

\CatchFileBetweenTags{\InvRM}{constants.tex}{InvRM}

\subsection{The Spacetime Partition (Weak Mixing Angle \texorpdfstring{$\sin^2 \theta_W$}{sin2thetaW})}
\label{subsec:spacetime_partition}
\label{subsec:weak_mixing}

\textbf{Architectural Requirement:} Under Rule 1 (Capacity Partitioning), a persistent system must continuously advance its state into the future. However, a local node is a heavy informational object. To update the node, the substrate must advance the total structural footprint of the system. It must strictly partition its operational bandwidth between the pure causal ``clock tick'' (Time) and the structural maintenance of the ``hardware'' (Space, Payload, and Wires). 

\textbf{The Geometric Primitives:} The total structural footprint of a node ($N_{sys}$) is the linear sum of its three orthogonal hardware components:

\begin{enumerate}
    \item \textbf{The Address Space ($R_M$):} The spacetime coordinate grid. The fundamental resonance ($\Delta=43$) distributed across the $D=4$ dimensions yields a manifold resolution of $R_M = 172$ coordinate slots.
    \item \textbf{The State Payload ($\nu \cdot \eta$):} The active chiral data ($\nu=16$) occupying the grid, strictly subject to the active transmission overhead ($\eta$), yielding $\approx 15.9$ active payload states.
    \item \textbf{The Interaction Bus ($\sigma$):} The $\sigma = 5$ independent generators required to physically connect the node to its neighbors.
\end{enumerate}

Summing these orthogonal components yields the Total System Footprint:
\begin{equation}
N_{sys} = \underbrace{R_M}_{\text{Address}}
+ \underbrace{\nu \cdot \eta}_{\text{State}}
+ \underbrace{\sigma}_{\text{Interaction}}
\approx \WeakAngleStepOneVal
\end{equation}

To maintain causal coherence, a portion of the bandwidth is allocated purely to the temporal update stream ($\Delta = 43$, the 1-dimensional causal period) against the total structural footprint of the system ($N_{sys}$):

\begin{equation}
\begin{split}
\text{Clock Ratio} &= \frac{\text{Temporal Updates}}{\text{Total System Footprint}} = \frac{\Delta}{N_{sys}} \\
&= \frac{43}{\WeakAngleStepOneVal} \approx \mathbf{\WeakAngleVal}
\end{split}
\end{equation}

\subsubsection{Physical Consequence: Electroweak Symmetry Breaking}

This geometric bandwidth ratio identically matches the Weak Mixing Angle ($\sin^2 \theta_W$) of the Standard Model. Rather than an arbitrary group-theoretic mixing parameter, it is the structural limit that dictates the physical bifurcation of the Electroweak force:

\begin{itemize}
    \item \textbf{The Photon ($\gamma$):} Represents the $\Delta$ numerator. Because it couples strictly to the temporal update stream, it carries no structural hardware load. It therefore propagates without geometric impedance as a massless carrier at the speed limit of the substrate ($c$).
    \item \textbf{The W/Z Bosons:} Represent the massive denominator ($N_{sys} - \Delta$). Because they must update the heavy spatial grid, state payload, and interaction bus, their propagation is severely impeded by the lattice structure, manifesting macroscopically as massive, short-range gauge bosons.
\end{itemize}

\vspace{0.5em}
\begin{tabular}{ll}
\textbf{Prediction:} & \num{\WeakAngleVal} \\
\textbf{Experiment (Direct Mass):} & \WeakAngleExperimentalValue \\
\textbf{Experiment (Global Fit):} & \WeakAngleGlobalExperimentalValue \\
\end{tabular}
\vspace{0.5em}

\textbf{The W-Boson Mass Tension:} The geometric prediction aligns with the \textbf{Direct Mass} measurement but deviates by \WeakAngleGlobalSigma$\sigma$ from the Standard Model Global Fit. The global fit derives $\sin^2\theta_W$ by propagating correlated constraints across the full electroweak parameter space using continuous perturbative regularization; the geometric prediction arises from a strictly finite, discrete capacity partition. These two procedures need not agree: if the substrate is genuinely discrete, continuous infinite-loop regularization would systematically shift the inferred mixing angle. Whether this \WeakAngleGlobalSigma$\sigma$ tension represents a structural prediction of the discrete substrate or a falsification of the framework depends on the definitive W-mass experimental determination. The direct PDG 2024 measurement ($M_W = 80.369 \pm 0.013\,\mathrm{GeV}$) is consistent with the geometric value; the SM global fit remains a genuine challenge this framework must resolve.

\textit{Note on Physical Fields:} The explicit mapping of this Spacetime Partition to the physical gauge fields---demonstrating why the temporal allocation ($\Delta$) manifests as the massless Photon ($\gamma$) and the spatial infrastructure allocation manifests as the massive $W/Z$ bosons---along with the derivation of the absolute link admittances ($g, g', g_3$), is detailed in \cref{app:gauge_admittance}.